{
\def\AMCbeginQuestion#1#2{\par\noindent{\bf (POSCOMP 2004 - Questão 02)}}
\begin{question}{} Seja $V$ um espaço vetorial real com produto interno. Para $x$ e $y$ vetores quaisquer de $V$, a igualdade $\|x+y\|=\|x\|+\|y\|$ é verdadeira se, e somente se,
\begin{choices}
\correctchoice{$x=0$, ou $y=0$, ou ($x\neq 0$ e $y=\lambda x$) onde $\lambda$ é um número real não-negativo.}
\wrongchoice{$x\neq 0$ e $y=\lambda x$ para todo número real $\lambda$.}
\wrongchoice{$x=0$, ou $y=0$.}
\wrongchoice{$x=0$, ou $y=0$, ou ($x\neq 0$ e $x,y$ são linearmente dependentes).}
\wrongchoice{$x=0$, ou $y=0$, ou ($x\neq 0$ e $x,y$ são linearmente independentes).}
\end{choices}
\end{question}
}
%
{
\def\AMCbeginQuestion#1#2{\par\noindent{\bf (POSCOMP 2004 - Questão 03)}}
\begin{question}{} Sobre a transformação linear $T:\mathrm{R}^{2}\to\mathrm{R}^{2}$ definida pela matriz $\left[\begin{array}{rr}1&0\\-1&0\end{array}\right]$ podemos dizer que
\begin{choices}
\correctchoice{a imagem é a reta $y=-x$ e o núcleo é a reta $x=0$}
\wrongchoice{a imagem é a reta $y=x$ e o núcleo é $\{(0,0)\}$}
\wrongchoice{a imagem é a reta $x=0$ e o núcleo é a reta $y=-x$}
\wrongchoice{a imagem é a reta $y=x$ e o núcleo é o $\mathrm{R}^{2}$}
\wrongchoice{a imagem é o $\mathrm{R}^{2}$ e o núcleo é a reta $y=x$}
\end{choices}
\end{question}
}
%
{
\def\AMCbeginQuestion#1#2{\par\noindent{\bf (POSCOMP 2004 - Questão 04)}}
\begin{question}{} A transformação $T(x,y)=\frac{1}{5}(-4x+3y,3x+4y)$ do plano no plano é
\begin{choices}
\correctchoice{uma reflexão através da reta $y=3x$}
\wrongchoice{uma expansão uniforme}
\wrongchoice{uma contração uniforme}
\wrongchoice{uma translação}
\wrongchoice{um cisalhamento horizontal}
\end{choices}
\end{question}
}
%
{
\def\AMCbeginQuestion#1#2{\par\noindent{\bf (POSCOMP 2004 - Questão 05)}}
\begin{question}{} No $\mathrm{R}^{3}$ com o produto escalar usual, tome $v=(1,-1,0)$ e o subespaço $S$ gerado por $\{(1,2,1),(-1,1,-1)\}$. O vetor de $S$ mais próximo de $v$ é
\begin{choices}
\correctchoice{$(1/2,-1,1/2)$}
\wrongchoice{$(1,-1,1)$}
\wrongchoice{$(2/3,-1,1/3)$}
\wrongchoice{$(1/100,-1,1/100)$}
\wrongchoice{$(2,-1,2)$}
\end{choices}
\end{question}
}
%
{
\def\AMCbeginQuestion#1#2{\par\noindent{\bf (POSCOMP 2004 - Questão 07)}}
\begin{question}{} Se $A$ é uma matriz $n\times n$ de entradas reais, cujas linhas são linearmente independentes, então não se pode afirmar que:
\begin{choices}
\correctchoice{$\det(A)=1$.}
\wrongchoice{$A$ é inversível.}
\wrongchoice{$A\cdot X=B$ tem solução única $X$ para todo $B\in\mathrm{R}^{n}$.}
\wrongchoice{As colunas de $A$ são linearmente independentes.}
\wrongchoice{O posto de $A$ é $n$.}
\end{choices}
\end{question}
}
%
{
\def\AMCbeginQuestion#1#2{\par\noindent{\bf (POSCOMP 2023 - Questão 01)}}
\begin{question}{} Utilize o método de Eliminação de Gauss para resolver o sistema a seguir:
\[
\begin{array}{rcl}
-3x+y+z&=&1\\
x-2y+z&=&4\\
-x+y-3z&=&-7
\end{array}
\]
\begin{choices}
\correctchoice{$(0,-1,2)$}
\wrongchoice{$(0,1,0)$}
\wrongchoice{$(2,0,2)$}
\wrongchoice{$(2,1,2)$}
\wrongchoice{$(-1,-1,-1)$}
\end{choices}
\end{question}
}
%
{
\def\AMCbeginQuestion#1#2{\par\noindent{\bf (POSCOMP 2023 - Questão 02)}}
\begin{question}{} Determine o valor de $x$ para que o vetor $(1,x,5)\in\mathrm{R}^{3}$ pertença ao subespaço $\langle(1,2,3),(1,1,1)\rangle$.
\begin{choices}
\correctchoice{$x=3$}
\wrongchoice{$x=0$}
\wrongchoice{$x=-1$}
\wrongchoice{$x=1$}
\wrongchoice{$x=7$}
\end{choices}
\end{question}
}
%
{
\def\AMCbeginQuestion#1#2{\par\noindent{\bf (POSCOMP 2023 - Questão 03)}}
\begin{question}{} Determine o polinômio característico de $J$:
\[
J=\left(\begin{array}{rrr}
2&-2&1\\
1&-1&1\\
1&-2&2
\end{array}\right)
\]
\begin{choices}
\correctchoice{$(1-x)^3$}
\wrongchoice{$0$}
\wrongchoice{$x$}
\wrongchoice{$(1-x)$}
\wrongchoice{$(1-x)^2$}
\end{choices}
\end{question}
}
